% Generated draft: requires mathematical and source review.
\documentclass[11pt]{article}
\usepackage[margin=25mm]{geometry}
\usepackage{fontspec}
\setmainfont{Latin Modern Roman}
\usepackage{amsmath,amssymb,mathtools,unicode-math}
\setmathfont{Latin Modern Math}
\usepackage{enumitem,needspace,xurl,hyperref}
\setlength{\parindent}{0pt}
\setlength{\parskip}{5pt}
\setlength{\emergencystretch}{3em}
\newcommand{\N}{\mathbb{N}}
\newcommand{\Z}{\mathbb{Z}}
\newcommand{\Q}{\mathbb{Q}}
\newcommand{\R}{\mathbb{R}}
\newcommand{\C}{\mathbb{C}}
\newcommand{\F}{\mathbb{F}}
\newcommand{\A}{\mathbb{A}}
\newcommand{\PP}{\mathbb{P}}
\newcommand{\E}{\mathbb{E}}
\newcommand{\eps}{\varepsilon}
\newcommand{\dd}{\,\mathrm{d}}
\begin{document}
\section*{UnsolvedMath ID: 30004607}
\textbf{Source problem code:} catalogue entry
\subsection{Definitions and mathematical statement}
Write \(M=A(R,G)^+\) for the elliptic Artin monoid named in the supplied record, and write \(X=S_{\mathrm{ell}}\setminus D_{\mathrm{ell}}\) for its elliptic discriminant complement. Fix the source's basepoint \(*\in X\). The group \(\pi_2(X,*)\) is the abelian group of based homotopy classes of based continuous maps \((S^2,s_0)\to(X,*)\).

The supplied record does not specify the admissible parameters \((R,G)\), a presentation of \(M\), the construction of \(S_{\mathrm{ell}}\) and \(D_{\mathrm{ell}}\), or the meaning of the relation-class and division notation below. Consequently, the following records the intended conjecture conditionally on those source constructions, rather than supplying a complete intrinsic definition of its proposed generators.

Let \(I\) denote the source-intended collection of pairs of noncancellative quadruples
\[
((a,b,c,d),(a',b,c,d')),
\]
with the displayed common entries \(b,c\). For each such pair, the source proposes a difference class
\[
\Delta_{a,a';b,c;d,d'}=
 a\backslash R(abd,acd)/d
 -a'\backslash R(a'bd',a'cd')/d'
 \quad\text{in }\pi_2(X,*).
\]
The intended conjecture is that these classes generate the second homotopy group:
\[
\pi_2(X,*)=
\left\langle
\Delta_{a,a';b,c;d,d'}:
((a,b,c,d),(a',b,c,d'))\in I
\right\rangle.
\]
Interpreting generation as generation of the underlying abelian group, this says that every \(\alpha\in\pi_2(X,*)\) is a finite integer linear combination of the displayed difference classes. The supplied record does not resolve whether an additional \(\pi_1(X,*)\)-module interpretation is intended. It also does not define noncancellative quadruples or establish how the displayed expressions yield well-defined based homotopy classes; those definitions and qualifications must be recovered from the cited source before this can be a fully specified mathematical statement.
\subsection{Short English statement}
The proposed conjecture asks whether differences associated with pairs of noncancellative quadruples sharing \(b,c\) generate \(\pi_2\) of the elliptic discriminant complement. The supplied record labels the problem open, but omits definitions needed to identify the proposed classes and their precise range.
\subsection{Sources}
\begin{itemize}
\item Local UnsolvedMath registry, supplied complete source record, permanent ID \(30004607\), exact source problem code \texttt{catalogue entry}, titled \emph{Homotopy Generators from Elliptic Artin Noncancellativity}; Oberwolfach Reports Open Problems collection. The record gives proposed year 2021 and names no individual proposer.
\item \emph{Logarithmic Vector Fields and Freeness of Divisors and Arrangements: New perspectives and applications} (2022), Oberwolfach Reports, pp.~225--300, DOI \texttt{10.4171/owr/2021/5}. Supplied report reference: \url{https://doi.org/10.4171/owr/2021/5}. The record describes Oberwolfach Reports as documenting mathematical workshops and their open-problem discussions.
\item Supplied literature source: \emph{Oberwolfach Report 5/2021}, \url{https://ems.press/content/serial-article-files/46883}. Literature assessment dated exactly \texttt{2026-08-22}: the elliptic Artin-monoid conjectures, including the proposed generators of the second homotopy group, remain conjectural, and the report cites the underlying work as in preparation. This is supplied dated curation derived from the OWR record, not an independently verified literature assessment; time-sensitive status and cited claims require independent verification.
\item The registry labels the statement \texttt{corrected\_verified} and reports that it was checked against the official Oberwolfach report, preserving its hypotheses, notation, and open target. These are supplied verification claims, not verification performed for this draft. The registry status is \texttt{open}.
\end{itemize}
\end{document}
